% Appendix C: Experiment 2 details. Design part written 2026-09-11 from
% code/synthetic/data_v2.py, code/experiments/exp2_intervention.py and
% results/exp2/calibration.json. Results tables are appended from
% results/exp2_summary.csv and results/exp2/PREDICTIONS.md when the grid completes.

\subsection{Generator}
Images are $16\times16$ on a torus with $S=256$ channels per pixel. Labels
$y_p\in\{-1,+1\}$ are i.i.d.\ Rademacher. Fixed orthonormal directions
$u,v_1,\dots,v_8\in\R^S$ are drawn once per seed (QR of a Gaussian matrix).
With $o_1,\dots,o_8$ the eight $3\times3$ offsets, the spectrum of pixel $p$
is
\[
 x_p=\alpha\,y_p\,u+\beta\sum_{d=1}^8\big(y_{p-o_d}+\tau\,\eta_{p,d}\big)v_d+\sigma\,\varepsilon_p,
 \qquad \eta_{p,d},\varepsilon_p\ \text{standard Gaussian},
\]
so that the pixel at $p+o_d$ carries $y_p$ along $v_d$: each pixel's label
is written into all eight neighbours, each along its own direction, while
$p$'s own $v$-coordinates carry the labels of $p$'s neighbours, which are
independent of $y_p$. Constants: $\alpha=1.645$, $\beta=5$, $\sigma=1$,
$\tau=1.7023$; $K=12$. The draw order is condition-independent, so the five
conditions generated from one seed share $y$, $\eta$ and $\varepsilon$:
\emph{iid}; \emph{reversed} ($y_{p-o_d}\to-y_{p-o_d}$); \emph{context-random}
(an independent label field replaces $y$ in the context term);
\emph{spectral-only} ($\beta=0$); \emph{context-only} ($\alpha=0$).
Training uses $256$ images; each test condition $64$ images; the spectral
probe is fitted on $64$ spectral-only images and evaluated on another $64$.
Seeds: directions $1000+s$, training data $2000+s$, test family $3000+s$,
probe $4000+s$ and $5000+s$, initialization $s$, for $s\in\{0,1,2\}$.

\subsection{Calibration and readiness (fixed before any training run)}
The oracle spectral accuracy is a Fisher discriminant on $u^\top x_p$; the
oracle contextual accuracy is a Fisher discriminant on the $3\times3$ window
of the eight projections $v_d^\top x$ (72 features), which contains all eight
copies of $y_p$. $\tau$ was set by bisection so that the contextual oracle
equals $0.950$ on seed-0 directions. Readiness is the same discriminant
through a random encoder $W\sim\mathcal N(0,1/S)^{12\times256}$: on the
encoder output of the pixel itself for spectral-only data, and on the
$3\times3$ window of encoder outputs for context-only data.

\begin{table}[h]
\centering\small
\caption{Experiment 2 calibration (\texttt{results/exp2/calibration.json}).
Oracle accuracies use the true directions; readiness uses a random encoder.}
\label{tab:exp2_calib}
\begin{tabular}{lccccc}
\toprule
seed & oracle spectral & oracle context & random enc.\ spectral & random enc.\ context & random enc.\ both \\
\midrule
0 & 0.951 & 0.950 & 0.647 & 0.910 & 0.915 \\
1 & 0.951 & 0.949 & 0.653 & 0.914 & 0.918 \\
2 & 0.948 & 0.950 & 0.608 & 0.905 & 0.906 \\
\bottomrule
\end{tabular}
\end{table}

The population optimum of the contextual cue alone is a cross-entropy of
about $0.126$ nats (spectral alone $0.127$); the threshold $0.15$ sits
just above that population optimum. A population optimum is not a bound
on the finite-sample training loss, which a wide head can drive lower by
fitting individual pixels, so training losses below $0.13$ are not by
themselves evidence of spectral learning; the probe and the shifted
accuracies are.

\subsection{Model, training and arms}
Encoder $W\in\R^{12\times256}$, no bias, $\mathcal N(0,1/S)$ initialization.
Head: $3\times3$ convolution to width $M$ with bias, ReLU, $3\times3$
convolution to one logit without bias, circular padding, PyTorch default
initialization; implemented as explicit matrix products with torus rolls
(identical to the convolution to $10^{-15}$ in double precision). Loss: mean
over the $65{,}536$ training pixels of $\log(1+e^{-y_pF_p})$. Full-batch
gradient descent, $\eta=10^{-3}$ for every parameter (the largest of
$\{10^{-4},3\cdot10^{-4},10^{-3},3\cdot10^{-3}\}$ with a monotone loss over
$300$ steps at $M=2048$), no momentum, clipping or decay; at most $40{,}000$
steps, stopping at loss $<0.10$. Snapshots at the first step with loss below
each of $0.6,0.5,0.4,0.3,0.2,0.15,0.10$; $L^\ast=0.30$ is the pre-registered
threshold. Gradient norms at a snapshot are those of the snapshotted
iterate. Arms: \texttt{sp} (standard parameterization) and \texttt{mup}
(readout $M^{-1/2}\gamma$, $\gamma=\sqrt M\times$default, same function at
initialization; under plain gradient descent equal to a readout-only rate
$1/M$) at $M\in\{2,4,8,32,128,512,2048\}$; \texttt{ctxfree} (training
context from an independent label field) and \texttt{frozen} (encoder fixed
at initialization) at $M\in\{2,8,128,512\}$; \texttt{lrmult} (readout rate
$\times\{1/16,1/4,1,4,16\}$) and \texttt{headlr} (whole-head rate
$\kappa\in\{1,1/16,1/256\}$, extended to $1/4096$ and $1/32768$ on seed 0
for the reversal outcome) at $M=32$. Three seeds.

\subsection{Pre-registration and amendments}
The plan (\texttt{paper/REFOCUS\_PLAN\_2026-09-10.md}, \S4) fixed the
generator, constants, arms, $L^\ast$, measures and predictions P1--P5 before
the grid. Amendments recorded before the corresponding runs: widths $2$ and
$4$ and fractional readout multipliers were added after a one-seed pilot at
$M\in\{8,2048\}$ showed complete suppression at $L^\ast$ (\S4.4a); the
whole-head arm, the exact $h_u$ export, paired test conditions and
same-state gradient norms were added after the design review (\S4.4b); the
$\kappa$ bracket was extended downward after the seed-0 pilot showed that
the probe responds to $\kappa$ while reversal accuracy does not within
$\{1,1/16,1/256\}$ (\S4.4c). Original predictions are reported as
registered; added widths and multipliers are identified as amendments.

\subsection{Results}
Tables generated by \texttt{paper/figures\_src/make\_tables\_exp2.py} from
\texttt{results/exp2\_summary.csv}; the registered verdicts are the
per-seed Spearman correlations listed at the end.
\IfFileExists{sections_iclr_v2/appendix_exp2_tables.tex}{\begin{table}[htbp]\centering\scriptsize
\caption{Experiment 2 at the pre-registered threshold $L^\ast=0.30$: mean$\pm$sd over seeds (number of seeds). Probe gain is the discriminant accuracy on the encoder output minus its initial value (per seed: 0.647, 0.658, 0.615, 0.659, 0.604); $h_u=\|P_{\mathrm{row}(W)}u\|^2$ (0.055 at initialization, seed 0).}\label{tab:exp2_star}
\begin{tabular}{llccccc}\toprule
arm & $M$ & probe gain & $h_u$ & reversed acc. & context-random acc. & spectral-only acc. \\ \midrule
\texttt{sp} & 2 & 0.043$\pm$0.024 (5) & 0.083$\pm$0.040 (5) & 0.133$\pm$0.007 (5) & 0.504$\pm$0.005 (5) & 0.526$\pm$0.016 (5) \\
\texttt{sp} & 4 & 0.020$\pm$0.012 (5) & 0.062$\pm$0.029 (5) & 0.127$\pm$0.006 (5) & 0.506$\pm$0.008 (5) & 0.526$\pm$0.017 (5) \\
\texttt{sp} & 8 & 0.024$\pm$0.012 (5) & 0.063$\pm$0.020 (5) & 0.123$\pm$0.002 (5) & 0.507$\pm$0.004 (5) & 0.526$\pm$0.013 (5) \\
\texttt{sp} & 32 & 0.016$\pm$0.010 (5) & 0.060$\pm$0.025 (5) & 0.118$\pm$0.008 (5) & 0.506$\pm$0.008 (5) & 0.524$\pm$0.020 (5) \\
\texttt{sp} & 128 & 0.011$\pm$0.005 (5) & 0.054$\pm$0.021 (5) & 0.116$\pm$0.010 (5) & 0.506$\pm$0.008 (5) & 0.523$\pm$0.016 (5) \\
\texttt{sp} & 512 & 0.006$\pm$0.003 (5) & 0.050$\pm$0.018 (5) & 0.115$\pm$0.009 (5) & 0.508$\pm$0.008 (5) & 0.525$\pm$0.018 (5) \\
\texttt{sp} & 2048 & 0.003$\pm$0.001 (5) & 0.048$\pm$0.017 (5) & 0.119$\pm$0.008 (5) & 0.508$\pm$0.008 (5) & 0.525$\pm$0.015 (5) \\
\midrule
\texttt{mup} & 2 & 0.056$\pm$0.030 (5) & 0.097$\pm$0.049 (5) & 0.134$\pm$0.008 (5) & 0.504$\pm$0.005 (5) & 0.531$\pm$0.020 (5) \\
\texttt{mup} & 4 & 0.032$\pm$0.017 (5) & 0.073$\pm$0.035 (5) & 0.126$\pm$0.008 (5) & 0.506$\pm$0.007 (5) & 0.529$\pm$0.020 (5) \\
\texttt{mup} & 8 & 0.039$\pm$0.021 (5) & 0.077$\pm$0.025 (5) & 0.119$\pm$0.002 (5) & 0.506$\pm$0.004 (5) & 0.531$\pm$0.018 (5) \\
\texttt{mup} & 32 & 0.034$\pm$0.017 (5) & 0.075$\pm$0.035 (5) & 0.112$\pm$0.007 (5) & 0.506$\pm$0.007 (5) & 0.530$\pm$0.025 (5) \\
\texttt{mup} & 128 & 0.026$\pm$0.012 (5) & 0.067$\pm$0.025 (5) & 0.108$\pm$0.009 (5) & 0.505$\pm$0.007 (5) & 0.524$\pm$0.013 (5) \\
\texttt{mup} & 512 & 0.036$\pm$0.021 (5) & 0.075$\pm$0.026 (5) & 0.109$\pm$0.007 (5) & 0.506$\pm$0.007 (5) & 0.529$\pm$0.016 (5) \\
\texttt{mup} & 2048 & 0.039$\pm$0.015 (5) & 0.077$\pm$0.026 (5) & 0.107$\pm$0.006 (5) & 0.506$\pm$0.007 (5) & 0.530$\pm$0.018 (5) \\
\midrule
\texttt{ctxfree} & 2 & 0.299$\pm$0.022 (5) & 0.858$\pm$0.035 (5) & 0.876$\pm$0.002 (5) & 0.876$\pm$0.002 (5) & 0.877$\pm$0.002 (5) \\
\texttt{ctxfree} & 8 & 0.286$\pm$0.025 (5) & 0.740$\pm$0.012 (5) & 0.892$\pm$0.004 (5) & 0.892$\pm$0.004 (5) & 0.897$\pm$0.004 (5) \\
\texttt{ctxfree} & 128 & 0.272$\pm$0.025 (5) & 0.651$\pm$0.010 (5) & 0.895$\pm$0.003 (5) & 0.894$\pm$0.002 (5) & 0.902$\pm$0.004 (5) \\
\texttt{ctxfree} & 512 & 0.256$\pm$0.025 (5) & 0.563$\pm$0.004 (5) & 0.883$\pm$0.004 (5) & 0.885$\pm$0.004 (5) & 0.888$\pm$0.004 (5) \\
\midrule
\texttt{frozen} & 2 & 0.000$\pm$0.000 (5) & 0.046$\pm$0.017 (5) & 0.140$\pm$0.009 (5) & 0.508$\pm$0.006 (5) & 0.523$\pm$0.016 (5) \\
\texttt{frozen} & 8 & 0.000$\pm$0.000 (5) & 0.046$\pm$0.017 (5) & 0.134$\pm$0.006 (5) & 0.510$\pm$0.007 (5) & 0.522$\pm$0.016 (5) \\
\texttt{frozen} & 128 & 0.000$\pm$0.000 (5) & 0.046$\pm$0.017 (5) & 0.127$\pm$0.008 (5) & 0.508$\pm$0.008 (5) & 0.524$\pm$0.017 (5) \\
\texttt{frozen} & 512 & 0.000$\pm$0.000 (5) & 0.046$\pm$0.017 (5) & 0.122$\pm$0.008 (5) & 0.508$\pm$0.009 (5) & 0.526$\pm$0.016 (5) \\
\bottomrule
\end{tabular}\end{table}
\begin{table}[htbp]\centering\scriptsize
\caption{Experiment 2 at the secondary threshold 0.15 (labelled; not pre-registered): mean$\pm$sd over seeds that reached it.}\label{tab:exp2_015}
\begin{tabular}{llcccc}\toprule
arm & $M$ & probe gain & $h_u$ & reversed acc. & context-random acc. \\ \midrule
\texttt{sp} & 2 & 0.213$\pm$0.030 (5) & 0.427$\pm$0.191 (5) & 0.117$\pm$0.069 (5) & 0.534$\pm$0.037 (5) \\
\texttt{sp} & 4 & 0.110$\pm$0.027 (5) & 0.167$\pm$0.049 (5) & 0.073$\pm$0.006 (5) & 0.507$\pm$0.006 (5) \\
\texttt{sp} & 8 & 0.120$\pm$0.019 (5) & 0.177$\pm$0.036 (5) & 0.073$\pm$0.005 (5) & 0.508$\pm$0.006 (5) \\
\texttt{sp} & 32 & 0.084$\pm$0.016 (5) & 0.130$\pm$0.048 (5) & 0.071$\pm$0.007 (5) & 0.508$\pm$0.008 (5) \\
\texttt{sp} & 128 & 0.061$\pm$0.016 (5) & 0.103$\pm$0.043 (5) & 0.070$\pm$0.007 (5) & 0.506$\pm$0.008 (5) \\
\texttt{sp} & 512 & 0.048$\pm$0.014 (5) & 0.088$\pm$0.034 (5) & 0.073$\pm$0.010 (5) & 0.508$\pm$0.008 (5) \\
\texttt{sp} & 2048 & 0.037$\pm$0.008 (5) & 0.075$\pm$0.026 (5) & 0.078$\pm$0.011 (5) & 0.509$\pm$0.009 (5) \\
\midrule
\texttt{mup} & 2 & 0.226$\pm$0.025 (5) & 0.470$\pm$0.188 (5) & 0.125$\pm$0.073 (5) & 0.538$\pm$0.041 (5) \\
\texttt{mup} & 4 & 0.160$\pm$0.031 (5) & 0.262$\pm$0.073 (5) & 0.079$\pm$0.009 (5) & 0.510$\pm$0.007 (5) \\
\texttt{mup} & 8 & 0.150$\pm$0.031 (5) & 0.236$\pm$0.064 (5) & 0.073$\pm$0.007 (5) & 0.509$\pm$0.006 (5) \\
\texttt{mup} & 32 & 0.112$\pm$0.032 (5) & 0.175$\pm$0.077 (5) & 0.068$\pm$0.007 (5) & 0.506$\pm$0.007 (5) \\
\texttt{mup} & 128 & 0.089$\pm$0.025 (5) & 0.141$\pm$0.051 (5) & 0.065$\pm$0.005 (5) & 0.504$\pm$0.006 (5) \\
\texttt{mup} & 512 & 0.117$\pm$0.036 (5) & 0.181$\pm$0.074 (5) & 0.067$\pm$0.006 (5) & 0.506$\pm$0.006 (5) \\
\texttt{mup} & 2048 & 0.118$\pm$0.025 (5) & 0.178$\pm$0.053 (5) & 0.067$\pm$0.005 (5) & 0.507$\pm$0.008 (5) \\
\midrule
\texttt{ctxfree} & 2 & 0.309$\pm$0.023 (3) & 0.990$\pm$0.000 (3) & 0.940$\pm$0.004 (3) & 0.940$\pm$0.005 (3) \\
\texttt{ctxfree} & 8 & 0.311$\pm$0.025 (5) & 0.976$\pm$0.004 (5) & 0.937$\pm$0.001 (5) & 0.938$\pm$0.002 (5) \\
\texttt{ctxfree} & 128 & 0.308$\pm$0.025 (5) & 0.937$\pm$0.001 (5) & 0.940$\pm$0.002 (5) & 0.940$\pm$0.001 (5) \\
\texttt{ctxfree} & 512 & 0.306$\pm$0.025 (5) & 0.913$\pm$0.003 (5) & 0.939$\pm$0.001 (5) & 0.939$\pm$0.001 (5) \\
\bottomrule
\end{tabular}\end{table}
\begin{table}[htbp]\centering\scriptsize
\caption{Experiment 2 at the secondary threshold 0.1 (labelled; not pre-registered): mean$\pm$sd over seeds that reached it.}\label{tab:exp2_010}
\begin{tabular}{llcccc}\toprule
arm & $M$ & probe gain & $h_u$ & reversed acc. & context-random acc. \\ \midrule
\texttt{sp} & 2 & 0.274$\pm$0.017 (4) & 0.637$\pm$0.054 (4) & 0.145$\pm$0.034 (4) & 0.558$\pm$0.020 (4) \\
\texttt{sp} & 4 & 0.234$\pm$0.021 (5) & 0.471$\pm$0.042 (5) & 0.093$\pm$0.008 (5) & 0.527$\pm$0.006 (5) \\
\texttt{sp} & 8 & 0.242$\pm$0.027 (5) & 0.500$\pm$0.028 (5) & 0.097$\pm$0.006 (5) & 0.530$\pm$0.006 (5) \\
\texttt{sp} & 32 & 0.230$\pm$0.020 (5) & 0.454$\pm$0.034 (5) & 0.096$\pm$0.008 (5) & 0.529$\pm$0.007 (5) \\
\texttt{sp} & 128 & 0.208$\pm$0.017 (5) & 0.382$\pm$0.033 (5) & 0.095$\pm$0.008 (5) & 0.527$\pm$0.008 (5) \\
\texttt{sp} & 512 & 0.174$\pm$0.017 (5) & 0.285$\pm$0.034 (5) & 0.096$\pm$0.008 (5) & 0.529$\pm$0.008 (5) \\
\texttt{sp} & 2048 & 0.147$\pm$0.016 (5) & 0.223$\pm$0.030 (5) & 0.107$\pm$0.012 (5) & 0.533$\pm$0.010 (5) \\
\midrule
\texttt{mup} & 2 & 0.276$\pm$0.018 (3) & 0.704$\pm$0.062 (3) & 0.173$\pm$0.061 (3) & 0.574$\pm$0.032 (3) \\
\texttt{mup} & 4 & 0.259$\pm$0.020 (5) & 0.586$\pm$0.056 (5) & 0.106$\pm$0.017 (5) & 0.534$\pm$0.010 (5) \\
\texttt{mup} & 8 & 0.256$\pm$0.024 (4) & 0.595$\pm$0.041 (4) & 0.097$\pm$0.010 (4) & 0.531$\pm$0.008 (4) \\
\texttt{mup} & 32 & 0.243$\pm$0.002 (2) & 0.591$\pm$0.024 (2) & 0.089$\pm$0.000 (2) & 0.528$\pm$0.001 (2) \\
\texttt{mup} & 128 & 0.237 (1) & 0.579 (1) & 0.084 (1) & 0.524 (1) \\
\texttt{mup} & 512 & 0.250$\pm$0.027 (3) & 0.559$\pm$0.045 (3) & 0.079$\pm$0.007 (3) & 0.521$\pm$0.006 (3) \\
\texttt{mup} & 2048 & 0.254$\pm$0.032 (2) & 0.540$\pm$0.016 (2) & 0.077$\pm$0.007 (2) & 0.519$\pm$0.010 (2) \\
\bottomrule
\end{tabular}\end{table}
\begin{table}[htbp]\centering\scriptsize
\caption{Whole-head rate multiplier $\kappa$ at $M=32$ (arm \texttt{headlr}) at $L^\ast$; rows marked $\dagger$ never reached $L^\ast$ within 40,000 steps and show the final state (loss in parentheses).}\label{tab:exp2_kappa}
\begin{tabular}{llrccccc}\toprule seed & $\kappa$ & step & probe gain & $h_u$ & reversed & context-random & spectral-only \\ \midrule
0 & 3.05176e-05$\dagger$ & 40000 (0.347) & 0.267 & 0.701 & 0.208 & 0.528 & 0.689 \\
0 & 0.000244141$\dagger$ & 40000 (0.320) & 0.266 & 0.693 & 0.184 & 0.529 & 0.697 \\
0 & 0.00390625 & 22708 & 0.223 & 0.475 & 0.152 & 0.519 & 0.641 \\
0 & 0.0625 & 7611 & 0.106 & 0.177 & 0.129 & 0.513 & 0.568 \\
0 & 1 & 1671 & 0.028 & 0.079 & 0.128 & 0.515 & 0.536 \\
1 & 0.00390625 & 34174 & 0.203 & 0.443 & 0.146 & 0.510 & 0.615 \\
1 & 0.0625 & 8977 & 0.100 & 0.181 & 0.129 & 0.509 & 0.575 \\
1 & 1 & 1580 & 0.025 & 0.085 & 0.121 & 0.511 & 0.551 \\
2 & 0.00390625 & 32068 & 0.205 & 0.301 & 0.130 & 0.500 & 0.570 \\
2 & 0.0625 & 8704 & 0.056 & 0.067 & 0.110 & 0.498 & 0.520 \\
2 & 1 & 1881 & 0.014 & 0.036 & 0.112 & 0.500 & 0.507 \\
3 & 0.00390625 & 39322 & 0.225 & 0.524 & 0.155 & 0.522 & 0.634 \\
3 & 0.0625 & 10120 & 0.046 & 0.104 & 0.127 & 0.511 & 0.536 \\
3 & 1 & 1731 & 0.008 & 0.069 & 0.118 & 0.510 & 0.525 \\
4 & 0.00390625 & 37418 & 0.221 & 0.316 & 0.133 & 0.502 & 0.561 \\
4 & 0.0625 & 9818 & 0.038 & 0.048 & 0.119 & 0.497 & 0.506 \\
4 & 1 & 1894 & 0.005 & 0.030 & 0.109 & 0.497 & 0.502 \\
\bottomrule\end{tabular}\end{table}
\begin{table}[htbp]\centering\scriptsize
\caption{Readout learning-rate multiplier at $M=32$ (arm \texttt{lrmult}) at $L^\ast$: mean$\pm$sd over 5 seeds.}\label{tab:exp2_lrmult}
\begin{tabular}{lcccc}\toprule multiplier & probe gain & $h_u$ & reversed acc. & context-random acc. \\ \midrule
0.0625 & 0.031$\pm$0.016 & 0.073$\pm$0.034 & 0.112$\pm$0.007 & 0.506$\pm$0.007 \\
0.25 & 0.023$\pm$0.015 & 0.067$\pm$0.030 & 0.114$\pm$0.007 & 0.506$\pm$0.007 \\
1 & 0.016$\pm$0.010 & 0.060$\pm$0.025 & 0.118$\pm$0.008 & 0.506$\pm$0.008 \\
4 & 0.011$\pm$0.005 & 0.055$\pm$0.021 & 0.124$\pm$0.008 & 0.507$\pm$0.007 \\
16 & 0.007$\pm$0.002 & 0.051$\pm$0.019 & 0.130$\pm$0.009 & 0.507$\pm$0.007 \\
\bottomrule\end{tabular}\end{table}
\begin{table}[htbp]\centering\scriptsize
\caption{Recovery at $M=32$: a fresh head with paired initialization trained for 20,000 steps on newly sampled context-random images with the encoder frozen at its $L^\ast$ checkpoint (\texttt{init} = the random initial encoder). Original = the run's own classifier at $L^\ast$; retrained = the new head; accuracies on the same paired test family.}\label{tab:exp2_recovery}
\begin{tabular}{llccccccc}\toprule seed & encoder & probe & retrain loss & orig.\ reversed & orig.\ ctx-random & retr.\ iid & retr.\ reversed & retr.\ ctx-random \\ \midrule
0 & \texttt{init} & 0.647 & 0.658 & -- & -- & 0.592 & 0.599 & 0.596 \\
0 & \texttt{kappa1} & 0.675 & 0.645 & 0.128 & 0.515 & 0.608 & 0.618 & 0.612 \\
0 & \texttt{kappa256} & 0.870 & 0.328 & 0.152 & 0.519 & 0.845 & 0.840 & 0.844 \\
1 & \texttt{init} & 0.658 & 0.659 & -- & -- & 0.596 & 0.590 & 0.593 \\
1 & \texttt{kappa1} & 0.683 & 0.658 & 0.121 & 0.511 & 0.594 & 0.591 & 0.592 \\
1 & \texttt{kappa256} & 0.861 & 0.447 & 0.146 & 0.510 & 0.776 & 0.771 & 0.773 \\
2 & \texttt{init} & 0.615 & 0.682 & -- & -- & 0.562 & 0.533 & 0.546 \\
2 & \texttt{kappa1} & 0.628 & 0.676 & 0.112 & 0.500 & 0.574 & 0.539 & 0.554 \\
2 & \texttt{kappa256} & 0.820 & 0.470 & 0.130 & 0.500 & 0.757 & 0.750 & 0.755 \\
3 & \texttt{init} & 0.659 & 0.647 & -- & -- & 0.614 & 0.605 & 0.612 \\
3 & \texttt{kappa1} & 0.667 & 0.653 & 0.118 & 0.510 & 0.604 & 0.600 & 0.603 \\
3 & \texttt{kappa256} & 0.884 & 0.378 & 0.155 & 0.522 & 0.817 & 0.808 & 0.811 \\
4 & \texttt{init} & 0.604 & 0.677 & -- & -- & 0.551 & 0.561 & 0.558 \\
4 & \texttt{kappa1} & 0.609 & 0.678 & 0.109 & 0.497 & 0.534 & 0.567 & 0.551 \\
4 & \texttt{kappa256} & 0.824 & 0.453 & 0.133 & 0.502 & 0.766 & 0.774 & 0.769 \\
\bottomrule\end{tabular}\end{table}
\paragraph{Runs that did not reach $L^\ast$.}
The following runs ended at the step budget above $L^\ast$: \texttt{headlr} $M=32$, $\kappa=0.000244141$, seed 0 (final loss 0.320 at step 40000); \texttt{headlr} $M=32$, $\kappa=3.05176e-05$, seed 0 (final loss 0.347 at step 40000).
\begin{table}[htbp]\centering\scriptsize
\caption{Registered predictions (REFOCUS\_PLAN \S4.4) evaluated at $L^\ast=0.30$ on the original five widths and on the amended grid; $\rho$ = Spearman rank correlation across widths (or multipliers), one value per registered seed (0--2; added seeds are reported in the mean$\pm$sd tables only). P2's second clause compares $a_u$ at $M=2048$.}\label{tab:exp2_verdicts}
\begin{tabular}{p{4.6cm}p{2.6cm}p{3.8cm}p{2.0cm}}\toprule criterion & grid & per-seed values & verdict \\ \midrule
P1: standard param., $\rho(a_u,M)\le-0.8$ and $\rho(\text{rev},M)\le-0.8$ in every seed & original & $a_u$: -0.70, -1.00, -1.00; rev: -0.60, -0.90, -0.60 & not met (1/3 seeds) \\
 & amended & $a_u$: -0.54, -1.00, -0.96; rev: -0.64, -0.96, -0.82 & not met \\
P2a: normalized readout, $|\rho(\cdot,M)|<0.5$ in at least two seeds & original & $a_u$: +0.10, -0.90, -0.30; probe: -0.20, -0.90, -0.30; rev: -0.90, -0.70, -0.70 & alignment met; probe met; reversal not met \\
 & amended & $a_u$: +0.25, -0.96, -0.54; rev: -0.89, -0.89, -0.89 & alignment not met; reversal not met \\
P2b: $a_u(2048)$ normalized $>$ standard in every seed & paired seeds & 0.0059 vs 0.0045; 0.0062 vs 0.0045; 0.0043 vs 0.0027 & met \\
P3: uninformative context, $a_u>0.25$ at every width and $|\rho(a_u,M)|<0.5$ & amended control widths $\{2,8,128,512\}$ & $a_u$ range 0.058--0.174; $\rho(a_u,M)$: -1.00, -1.00, -1.00; $\rho(\text{probe},M)$: -1.00, -1.00, -1.00 & cutoff not met; flatness not met (probe 0.936$\to$0.893) \\
P4: readout multiplier, $a_u$ and reversal accuracy decrease with the multiplier in every seed & run multipliers $\{1/16,1/4,1,4,16\}$ (64 dropped) & $a_u$: -1.00, -1.00, -1.00; rev: +1.00, +1.00, +1.00 & $a_u$ met; reversal not met (opposite sign) \\
P5: frozen encoder, reversal accuracy $<0.5$ at every width & amended widths & max 0.146 & met \\
P6 (amendment 4.4b): whole-head $\kappa$, probe, $h_u$ and reversal accuracy increase as $\kappa$ decreases, $\rho\le-0.8$ in every seed & $\kappa\in\{1,1/16,1/256\}$, three seeds & probe: -1.00, -1.00, -1.00; $h_u$: -1.00, -1.00, -1.00; rev: -1.00, -1.00, -0.50 & probe and $h_u$ met; reversal not met (2/3 seeds) \\
\bottomrule\end{tabular}\end{table}
}{\emph{[tables pending]}}
